\documentclass[10pt,letterpaper]{article}
\usepackage[T1]{fontenc}
\usepackage[utf8]{inputenc}
\usepackage[spanish,es-tabla]{babel}
\usepackage{csquotes}
\usepackage{mathpazo}
\usepackage[scaled=0.9]{helvet}
\usepackage{courier}
\usepackage{calc}
\usepackage[top=2cm, bottom=2cm, left=2cm, right=2cm]{geometry}
\usepackage{xcolor}
\usepackage{booktabs}
\usepackage{longtable}
\usepackage{array}
\usepackage{ragged2e}
\usepackage{xurl}
\usepackage[strings]{underscore}
\usepackage{fancyhdr}
\usepackage{sectsty}
\usepackage{tocloft}
\usepackage[colorlinks=true,linkcolor=blue,urlcolor=blue,citecolor=blue,breaklinks=true]{hyperref}
\usepackage{enumitem}

\definecolor{uncpblue}{HTML}{1A237E}
\allsectionsfont{\color{uncpblue}}

\renewcommand{\cftsecleader}{\cftdotfill{\cftsecdotsep}}
\cftpagenumbersoff{section}
\cftpagenumbersoff{subsection}
\cftpagenumbersoff{subsubsection}
\renewcommand{\cftsecfont}{\normalfont}
\renewcommand{\cftsubsecfont}{\normalfont}
\renewcommand{\cftsubsubsecfont}{\normalfont}
\setlength{\parindent}{0pt}
\setlength{\parskip}{4pt}
\setlength{\tabcolsep}{3pt}
\renewcommand{\arraystretch}{0.85}

\pagestyle{fancy}
\fancyhf{}
\fancyhead[L]{\small\textcolor{uncpblue}{\textbf{SGSI-UNCP}}}
\fancyhead[R]{\small\textcolor{uncpblue}{POL-SGSI-07}}
\fancyfoot[C]{\thepage}
\renewcommand{\headrulewidth}{0.4pt}
\renewcommand{\footrulewidth}{0.4pt}
\setlength{\headheight}{14pt}

\setlist{nosep,leftmargin=*,after=\vspace{2pt}}
\setlength{\emergencystretch}{2em}

\newcommand{\makecover}{
\begin{titlepage}
\centering
\vspace*{2cm}
{\Huge\bfseries\color{uncpblue} SGSI-UNCP\par}
\vspace{0.7cm}
{\LARGE Política de Seguridad Física y Áreas Seguras\par}
\vspace{0.5cm}
{\large Documento Base del Sistema de Gestión\\de Seguridad de la Información\par}
\vspace{1.5cm}
{\small Universidad Nacional del Centro del Perú\par}
\vfill
\end{titlepage}
}

\begin{document}

\makecover
\setcounter{tocdepth}{2}
\RaggedRight
\tableofcontents
\newpage

\RaggedRight
\section{Política de Seguridad Física y Áreas
Seguras}\label{poluxedtica-de-seguridad-fuxedsica-y-uxe1reas-seguras}

\textbf{Organización:} Universidad Nacional del Centro del Perú (UNCP)
\textbf{Aprobado por:} Comité de Gobierno y Transformación Digital
\textbf{Norma:} ISO/IEC 27001:2022 (Controles A.7.1--A.7.14)
\textbf{Alineamiento:} P-SGSI-03 (Gestión de Accesos), P-SGSI-07
(Eliminación Segura), F-SGSI-06 (Bitácora de Acceso a Áreas Críticas)

\begin{center}\rule{0.5\linewidth}{0.5pt}\end{center}

\subsection{Objetivo}\label{objetivo}

Establecer las medidas de seguridad física para proteger las
instalaciones, equipos y activos de información de la UNCP contra
accesos no autorizados, daños, interferencias y amenazas ambientales,
garantizando la continuidad de los servicios críticos.

\subsection{Alcance}\label{alcance}

Esta política aplica a todas las instalaciones de la UNCP, incluyendo
sus \textbf{4 sedes} (Huancayo, Mantaro, Satipo, Tarma), con énfasis en
las \textbf{áreas críticas} definidas en la sección 3. Aplica a todo el
personal, visitantes, contratistas y proveedores que ingresen a dichas
instalaciones.

\subsection{Clasificación de Áreas}\label{clasificaciuxf3n-de-uxe1reas}

{\setlength{\RaggedRightRightskip}{0pt plus 5em}\small\sloppy \begin{longtable}[]{@{}
  >{\hspace{0pt}\RaggedRight\arraybackslash}p{(\linewidth - 4\tabcolsep) * \real{0.3333}}
  >{\hspace{0pt}\RaggedRight\arraybackslash}p{(\linewidth - 4\tabcolsep) * \real{0.3333}}
  >{\hspace{0pt}\RaggedRight\arraybackslash}p{(\linewidth - 4\tabcolsep) * \real{0.3333}}@{}}
\toprule\noalign{}
\begin{minipage}[b]{\linewidth}\raggedright
Nivel
\end{minipage} & \begin{minipage}[b]{\linewidth}\raggedright
Área
\end{minipage} & \begin{minipage}[b]{\linewidth}\raggedright
Criterio
\end{minipage} \\
\midrule\noalign{}
\endhead
\bottomrule\noalign{}
\endlastfoot
\textbf{Crítica} & Datacenter principal, sala de servidores, cuarto de
comunicaciones, Archivo Central & Almacenan o procesan activos de
información críticos. Acceso restringido al mínimo personal
autorizado \\
\textbf{Restringida} & Oficina de OTI, Tesorería, Registros Académicos,
RRHH, Centro Médico, sala de equipos de red por facultad & Manejan
información confidencial o equipos sensibles. Acceso limitado al
personal del área \\
\textbf{Controlada} & Oficinas administrativas, salas de reuniones,
bibliotecas, laboratorios de cómputo & Acceso con identificación
institucional o registro de ingreso \\
\textbf{Pública} & Hall principal, auditorios, cafetería, áreas comunes
& Acceso libre al público con identificación en recepción \\
\end{longtable}\normalsize}

\subsection{Perímetros de Seguridad
Física}\label{peruxedmetros-de-seguridad-fuxedsica}

\subsubsection{4.1 Datacenter y Salas de Servidores (Área
Crítica)}\label{datacenter-y-salas-de-servidores-uxe1rea-cruxedtica}

\begin{itemize}
\tightlist
\item
  El Datacenter debe contar con \textbf{barreras físicas sólidas} (muros
  de concreto desde losa a losa, puerta blindada con cerradura
  multipunto).
\item
  Sistema de \textbf{control de acceso electrónico} con autenticación
  biométrica (huella dactilar) + tarjeta de proximidad, integrado a un
  sistema centralizado de registro.
\item
  \textbf{Videovigilancia (CCTV)} continua con cámaras en todos los
  accesos, interior de la sala y perímetro exterior. Las grabaciones
  deben conservarse por un mínimo de \textbf{90 días}.
\item
  \textbf{Sistema de detección de intrusiones} (sensores de puerta,
  detectores de movimiento) activo 24/7, monitoreado por seguridad
  patrimonial.
\item
  \textbf{Bitácora de acceso} digital o física (\textbf{F-SGSI-06}) que
  registre toda entrada y salida, incluyendo personal, visitantes y
  contratistas.
\end{itemize}

\subsubsection{4.2 Perímetro del Campus}\label{peruxedmetro-del-campus}

\begin{itemize}
\tightlist
\item
  El ingreso al campus principal debe contar con \textbf{control de
  acceso} (caseta de vigilancia, barrera vehicular, identificación de
  visitantes).
\item
  El perímetro debe estar delimitado por \textbf{cercos o muros
  perimetrales} en buen estado.
\item
  Las puertas de acceso peatonal y vehicular deben permanecer cerradas o
  controladas durante el horario no laboral.
\end{itemize}

\subsection{Control de Acceso a Áreas
Críticas}\label{control-de-acceso-a-uxe1reas-cruxedticas}

\subsubsection{5.1 Personal Autorizado}\label{personal-autorizado}

\begin{itemize}
\tightlist
\item
  Solo el personal de la OTI con autorización expresa del Jefe de la OTI
  puede ingresar al Datacenter. La lista de autorizados debe revisarse
  trimestralmente.
\item
  El acceso a áreas Restringidas se limita al personal asignado a esa
  unidad. Personal de otras áreas solo puede ingresar acompañado por un
  miembro del área.
\end{itemize}

\subsubsection{5.2 Visitantes y
Contratistas}\label{visitantes-y-contratistas}

\begin{itemize}
\tightlist
\item
  Todo visitante debe:

  \begin{enumerate}
  \def\labelenumi{\arabic{enumi}.}
  \tightlist
  \item
    Registrarse en recepción (nombre, DNI, empresa, motivo, hora de
    ingreso y salida, persona a quien visita).
  \item
    Portar un \textbf{solapín o carnet de visitante} visible en todo
    momento.
  \item
    Ser acompañado por un empleado autorizado de la UNCP durante su
    permanencia en áreas Críticas o Restringidas.
  \end{enumerate}
\item
  Los contratistas de mantenimiento deben presentar orden de servicio y
  registrar sus herramientas y equipos al ingreso.
\end{itemize}

\subsubsection{5.3 Acceso Fuera del Horario
Laboral}\label{acceso-fuera-del-horario-laboral}

\begin{itemize}
\tightlist
\item
  El acceso a las instalaciones fuera del horario laboral (19:00 -
  07:00) y fines de semana debe ser autorizado por el jefe inmediato y
  registrado en el sistema de control de acceso o bitácora.
\item
  El personal de seguridad patrimonial debe verificar la identidad y
  autorización antes de permitir el ingreso.
\end{itemize}

\subsection{Protección contra Amenazas
Ambientales}\label{protecciuxf3n-contra-amenazas-ambientales}

\subsubsection{6.1 Incendios}\label{incendios}

\begin{itemize}
\tightlist
\item
  El Datacenter debe contar con sistema de \textbf{detección temprana de
  incendios} (sensores de humo, temperatura, aspiración VESDA).
\item
  Sistema de \textbf{extinción por gas limpio} (FM-200, Novec 1230 o
  equivalente) que no dañe los equipos electrónicos. No se aceptan
  sistemas de agua en el Datacenter.
\item
  Extintores portátiles de CO2 en puntos estratégicos del Datacenter y
  pasillos.
\item
  Pruebas periódicas del sistema contra incendios según la frecuencia
  indicada por el fabricante o la normativa municipal, coordinadas con
  la OTI para minimizar falsas alarmas.
\end{itemize}

\subsubsection{6.2 Clima y Humedad}\label{clima-y-humedad}

\begin{itemize}
\tightlist
\item
  El Datacenter debe mantener condiciones ambientales controladas:

  \begin{itemize}
  \tightlist
  \item
    \textbf{Temperatura:} 18°C a 24°C (recomendado 22°C).
  \item
    \textbf{Humedad relativa:} 40\% a 60\%.
  \end{itemize}
\item
  Sistema de monitoreo ambiental con alarmas automáticas ante
  desviaciones.
\item
  Aire acondicionado de precisión (CRAC/CRAH) con redundancia N+1.
\end{itemize}

\subsubsection{6.3 Energía Eléctrica}\label{energuxeda-eluxe9ctrica}

\begin{itemize}
\tightlist
\item
  \textbf{UPS (Sistema de Alimentación Ininterrumpida)} con autonomía
  suficiente para un apagado controlado de todos los servidores (mínimo
  30 minutos a carga completa).
\item
  \textbf{Generador eléctrico de respaldo} con capacidad para mantener
  todo el Datacenter operativo por al menos 24 horas.
\item
  Pruebas mensuales del generador bajo carga.
\item
  Cableado eléctrico y de datos segregado (canaletas separadas) para
  evitar interferencias electromagnéticas.
\end{itemize}

\subsection{Seguridad de las Áreas de Trabajo (Escritorio y Pantalla
Limpia)}\label{seguridad-de-las-uxe1reas-de-trabajo-escritorio-y-pantalla-limpia}

\begin{itemize}
\tightlist
\item
  Todo el personal debe mantener su escritorio \textbf{limpio} al
  finalizar la jornada laboral:

  \begin{itemize}
  \tightlist
  \item
    Documentos clasificados como Confidenciales o de Uso Interno
    guardados bajo llave o en archivadores cerrados.
  \item
    Dispositivos extraíbles (USB, discos externos) retirados y
    guardados.
  \end{itemize}
\item
  Las pantallas de computadora deben bloquearse automáticamente tras
  \textbf{5 minutos} de inactividad (configuración forzada por política
  de grupo/MDM).
\item
  No se debe dejar información sensible visible en la pantalla cuando el
  puesto de trabajo está desatendido.
\end{itemize}

\subsection{Seguridad de Activos Fuera de las
Instalaciones}\label{seguridad-de-activos-fuera-de-las-instalaciones}

\begin{itemize}
\tightlist
\item
  Los equipos portátiles (laptops, tablets) que sean retirados del
  campus deben:

  \begin{itemize}
  \tightlist
  \item
    Tener el cifrado obligatorio activo (según POL-SGSI-05).
  \item
    Ser transportados en fundas o maletines que no identifiquen su
    contenido.
  \item
    No dejarse desatendidos en vehículos o lugares públicos.
  \end{itemize}
\item
  La retirada temporal de equipos del campus debe ser registrada en la
  OTI mediante el formato \textbf{F-SGSI-03}, indicando destino,
  responsable y fecha estimada de retorno.
\end{itemize}

\subsection{Mantenimiento de Equipos}\label{mantenimiento-de-equipos}

\begin{itemize}
\tightlist
\item
  Todo mantenimiento de equipos del Datacenter o áreas críticas debe:

  \begin{itemize}
  \tightlist
  \item
    Realizarse dentro de una \textbf{ventana de mantenimiento}
    previamente aprobada por el Jefe de la OTI.
  \item
    Ser ejecutado por personal autorizado (interno o del fabricante).
  \item
    Quedar registrado en una bitácora de mantenimiento (fecha, equipo
    intervenido, técnico, trabajo realizado, duración).
  \end{itemize}
\item
  El personal de mantenimiento externo debe estar acompañado en todo
  momento por un miembro de la OTI.
\end{itemize}

\subsection{Cableado de Seguridad}\label{cableado-de-seguridad}

\begin{itemize}
\tightlist
\item
  El cableado de red y eléctrico debe estar protegido contra
  interceptaciones y daños:

  \begin{itemize}
  \tightlist
  \item
    Cableado estructurado en canaletas cerradas o tuberías.
  \item
    Los paneles de parcheo (patch panels) deben estar en racks cerrados
    con llave o en cuartos de comunicaciones con acceso restringido.
  \item
    El cableado de fibra óptica entre edificios debe estar canalizado
    subterráneamente o protegido mecánicamente.
  \end{itemize}
\end{itemize}

\subsection{Eliminación Segura de Activos
Físicos}\label{eliminaciuxf3n-segura-de-activos-fuxedsicos}

La eliminación o retiro definitivo de activos físicos (equipos,
mobiliario que contenía información) debe realizarse siguiendo el
\textbf{P-SGSI-07 (Procedimiento de Eliminación Segura)} y registrarse
en el \textbf{F-SGSI-05 (Acta de Eliminación Segura)}.

\subsection{Responsabilidades}\label{responsabilidades}

{\setlength{\RaggedRightRightskip}{0pt plus 5em}\small\sloppy \begin{longtable}[]{@{}
  >{\hspace{0pt}\RaggedRight\arraybackslash}p{(\linewidth - 2\tabcolsep) * \real{0.5000}}
  >{\hspace{0pt}\RaggedRight\arraybackslash}p{(\linewidth - 2\tabcolsep) * \real{0.5000}}@{}}
\toprule\noalign{}
\begin{minipage}[b]{\linewidth}\raggedright
Rol
\end{minipage} & \begin{minipage}[b]{\linewidth}\raggedright
Responsabilidad
\end{minipage} \\
\midrule\noalign{}
\endhead
\bottomrule\noalign{}
\endlastfoot
\textbf{Oficina de Seguridad Patrimonial} & Controlar el acceso físico
al campus; gestionar la vigilancia perimetral y CCTV; mantener las
bitácoras de visitantes. \\
\textbf{OTI} & Gestionar el control de acceso electrónico del
Datacenter; monitorear las condiciones ambientales; ejecutar el
mantenimiento de equipos. \\
\textbf{Oficial de Seguridad} & Supervisar el cumplimiento de esta
política; revisar trimestralmente las listas de acceso autorizado. \\
\textbf{Todo el Personal} & Cumplir las normas de escritorio limpio,
bloqueo de pantalla y registro de visitantes. \\
\end{longtable}\normalsize}

\subsection{Monitoreo y Métricas}\label{monitoreo-y-muxe9tricas}

{\setlength{\RaggedRightRightskip}{0pt plus 5em}\small\sloppy \begin{longtable}[]{@{}
  >{\hspace{0pt}\RaggedRight\arraybackslash}p{(\linewidth - 6\tabcolsep) * \real{0.2500}}
  >{\hspace{0pt}\RaggedRight\arraybackslash}p{(\linewidth - 6\tabcolsep) * \real{0.2500}}
  >{\hspace{0pt}\RaggedRight\arraybackslash}p{(\linewidth - 6\tabcolsep) * \real{0.2500}}
  >{\hspace{0pt}\RaggedRight\arraybackslash}p{(\linewidth - 6\tabcolsep) * \real{0.2500}}@{}}
\toprule\noalign{}
\begin{minipage}[b]{\linewidth}\raggedright
Indicador
\end{minipage} & \begin{minipage}[b]{\linewidth}\raggedright
Meta
\end{minipage} & \begin{minipage}[b]{\linewidth}\raggedright
Frecuencia
\end{minipage} & \begin{minipage}[b]{\linewidth}\raggedright
Fuente
\end{minipage} \\
\midrule\noalign{}
\endhead
\bottomrule\noalign{}
\endlastfoot
Accesos no autorizados al Datacenter & 0 & Mensual & Bitácora de
acceso \\
Disponibilidad del UPS/generador & 100\% & Mensual & Registro de
pruebas \\
Temperatura del Datacenter fuera de rango & \textless{} 1 hora
acumulada/mes & Continuo & Monitoreo ambiental \\
Cumplimiento de escritorio limpio & \textgreater{} 90\% & Trimestral &
Auditoría visual aleatoria \\
Tiempo de retención de CCTV & \textgreater= 90 días & Mensual &
Verificación del sistema \\
\end{longtable}\normalsize}

\subsection{Documentos Relacionados}\label{documentos-relacionados}

{\setlength{\RaggedRightRightskip}{0pt plus 5em}\small\sloppy \begin{longtable}[]{@{}
  >{\hspace{0pt}\RaggedRight\arraybackslash}p{(\linewidth - 2\tabcolsep) * \real{0.5000}}
  >{\hspace{0pt}\RaggedRight\arraybackslash}p{(\linewidth - 2\tabcolsep) * \real{0.5000}}@{}}
\toprule\noalign{}
\begin{minipage}[b]{\linewidth}\raggedright
Código
\end{minipage} & \begin{minipage}[b]{\linewidth}\raggedright
Nombre
\end{minipage} \\
\midrule\noalign{}
\endhead
\bottomrule\noalign{}
\endlastfoot
P-SGSI-03 & Gestión de Identidades y Control de Acceso \\
P-SGSI-05 & Resiliencia y Continuidad en Nube Híbrida \\
P-SGSI-07 & Procedimiento de Eliminación Segura de Información \\
POL-SGSI-05 & Política de Seguridad para Dispositivos Móviles del
Personal \\
F-SGSI-05 & Acta de Eliminación Segura de Activos \\
F-SGSI-06 & Bitácora de Acceso a Áreas Críticas \\
\end{longtable}\normalsize}

\subsection{***}\label{section}

\end{document}
